\section{A finite-band positive-tilt operation}\label{sec:tilt}

The following lemma isolates the additional obligation between positive mean
drift and a true escape estimate. It concerns a normal completely positive bath
channel. No classical trajectory representation or intermediate measurement is
assumed.

\begin{lemma}[Bounded corrected drift gives exponential coordinate escape]
\label{lem:tilt}
Let $\mathcal E=\bigoplus_{\ell\ge0}\mathcal E_\ell$, with each level finite
dimensional, and let $L=\sum_{\ell\ge0}\ell P_\ell$. Let $\Phi$ be a normal
trace-preserving channel with an isometric dilation
$W:\mathcal E\to\mathcal E\otimes\mathcal K$ of depth bandwidth one:
\[
 (P_j\otimes I)WP_\ell=0\quad\hbox{if }|j-\ell|>1.
\]
Suppose $C=C^\dagger$ is bounded, $\|C\|\le M$, and on the finite-depth core
\[
 K=\Phi^*(L)-L+\Phi^*(C)-C
\]
extends to a bounded operator satisfying $K\ge\nu I$, with $\nu>0$.
For any $t$ obeying
\[
 0<t\le1,\qquad tM\le\tfrac12,\qquad
 t(5+8M)\le\tfrac{\nu}{2},\qquad \nu t<3,
 \tag{T1}
\]
define
\[
 A_t=I-tC,\qquad X_t=e^{-tL/2}A_te^{-tL/2},
 \qquad \lambda=1-\nu t/3.
\]
Then $X_t$ is a positive, injective compact operator,
\[
 (1-tM)e^{-tL}\le X_t\le(1+tM)e^{-tL},
 \qquad \Phi^*(X_t)\le\lambda X_t,\quad 0<\lambda<1.
 \tag{T2}
\]
For every normal density $\tau$, every $n\ge0$, and every $h\ge0$,
\[
 \Tr[P_{\le h}\Phi^n(\tau)]
 \le \min\!\left\{1,\,
 \frac{1+tM}{1-tM}e^{th}\lambda^n\right\}
 \le \min\{1,\,3e^{th-\nu tn/3}\}.
 \tag{T3}
\]
In particular, for every $v<\nu/3$ the mass at depths at most $vn$
decays exponentially. There is no invariant normal state, and every normal
orbit converges to zero against compact bath observables. The constants may
depend on the fixed channel and on $C$, but the displayed estimate is uniform
over initial normal densities.
\end{lemma}

\begin{proof}
Write $W=W_{-1}+W_0+W_1$ in its grade components. Each has norm at most one:
it is the corresponding Fourier coefficient of
$(e^{iuL}\otimes I)We^{-iuL}$, equivalently the strong integral of contractions
against a scalar phase. The finite-band hypothesis makes the sum exact.
The weighted dilation has the bounded extension
\[
 W_t=(e^{-tL/2}\otimes I)We^{tL/2}
 =e^{t/2}W_{-1}+W_0+e^{-t/2}W_1.
\]
This identity initially means equality on finite-depth vectors; the finite
sum defines the bounded extension, without applying an unbounded conjugation
to an arbitrary vector.

Put $D_t=W_t-W$ and $D_1'=(W_{-1}-W_1)/2$. For $0\le t\le1$,
\[
 \|D_t\|\le2\sinh(t/2)\le2t,\qquad
 \|D_t-tD_1'\|\le2(\cosh(t/2)-1)\le t^2/2.
 \tag{T4}
\]
Both scalar inequalities follow from their nonnegative power series on
$0\le t/2\le1/2$. Consequently
\[
 \|W_t^\dagger W_t-I-t(W^\dagger D_1'+D_1'{}^\dagger W)\|
 \le5t^2.
 \tag{T5}
\]
The derivative is $W^\dagger D_1'+D_1'{}^\dagger W
=L-\Phi^*(L)$ on the core and hence as a bounded operator. Define the bounded
positive tilted map
\[
 \mathfrak T_t(B)=W_t^\dagger(B\otimes I)W_t.
\]
For bounded $C$,
\[
 \|\mathfrak T_t(C)-\Phi^*(C)\|
 \le M(2\|D_t\|+\|D_t\|^2)\le8Mt.
 \tag{T6}
\]
Combining (T5)--(T6) yields the actual operator Taylor bound
\[
 \mathfrak T_t(A_t)-A_t=-tK+\mathcal R_t,\qquad
 \|\mathcal R_t\|\le(5+8M)t^2.
 \tag{T7}
\]
Under (T1), this is at most $-\nu tI/2$. Since $A_t\le3I/2$,
\[
 \mathfrak T_t(A_t)\le A_t-\nu tI/2
 \le(1-\nu t/3)A_t.
\]
Conjugate this bounded inequality by $e^{-tL/2}$. The finite-band
factorization gives
$\Phi^*(X_t)=e^{-tL/2}\mathfrak T_t(A_t)e^{-tL/2}$, first as forms on the
core and then everywhere by boundedness. This proves (T2). In particular,
we have not exponentiated an operator inequality for $L+C$.

Iteration and normal duality give
$\Tr[X_t\Phi^n(\tau)]\le\lambda^n\Tr(X_t\tau)\le(1+tM)\lambda^n$.
Also $P_{\le h}\le e^{th}e^{-tL}$. The lower bound in (T2) therefore proves
(T3), including its uniformity over $\tau$.

Injectivity follows from $A_t\ge I/2$ and injectivity of $e^{-tL/2}$.
Compactness follows from finite-dimensional levels and the decay of
$e^{-t\ell/2}$. If a normal invariant density existed, pairing (T2) with it
would force its strictly positive $X_t$ expectation to vanish. For each
fixed $h$, (T3) gives $\Tr(P_{\le h}\Phi^n(\tau))\to0$. Every finite-rank
operator is approximable in norm by its finite-coordinate compression;
uniform trace norm one then extends this convergence to every compact
operator. This is mass escape, not trace-norm convergence to the zero
operator.
\end{proof}

\paragraph{What is being reused.}
A receiver of this operation must supply the actual normal channel, a finite
depth-bandwidth dilation, a bounded correction, and a positive operator lower
bound for its corrected drift. Mere growth of a mean, or an uncontrolled
boundary term in a truncated drift calculation, is not a substitute. For
bandwidth greater than one the same derivation changes the explicit Taylor
constants; that extension is not needed here. Infinite-dimensional normality
is used in taking expectations and distinguishing an escaped probability mass
from a normal invariant state.

\paragraph{Established method, specific new verification.}
Exponential position conjugation, positive superharmonic operators and
Lyapunov criteria have primary precedents in quantum random-walk and
birth-and-death theory. The point of this lemma is to provide the exact
finite-band, rooted coherent-channel interface needed here, including a
uniform second-order remainder and the full operator positivity. It does
not assert a new generic principle of exponential tilting.
